\begin{abstract}

\vspace{-0.4em}

Foundation models are transforming business workflows and boosting productivity, yet they remain noticeably absent from some engineering domains such as power system analysis, where strict physical consistency must be enforced.\\

We present \genco (GEometric Neural Corrective Optimizer), a unified neural solver for steady-state transmission grid analysis that handles power flow (PF), optimal power flow (OPF), and state estimation (SE) within a single architecture, eliminating task-specific pipelines by operating on a shared grid representation. To support advances in neural power system solvers, we introduce the open-source GridFM Development Framework (\cref{fig:motivational_example}) that standardizes synthetic data generation and training in a low-code environment. We also release large-scale datasets with millions of PF and OPF scenarios across diverse grid topologies to further support reproducible benchmarking.\\

We evaluate \genco on the latest benchmark datasets (\pfdelta, OPFData), against all current state-of-the-art neural solvers, as well as classical solvers (Newton–Raphson, IPOPT) and real-world Hydro-Québec SCADA data. For large-scale PF, \genco recovers the full AC operating state---including voltage magnitudes and reactive power that DC-PF cannot provide---while matching DC-PF-level active power-balance residuals, with up to \maxpfspeedups speedups over Newton–Raphson at only \dcpfspeedupovergenco the runtime of DC-PF. For OPF, it achieves up to \maxopfspeedups speedups over IPOPT while improving feasibility, optimality, and runtime over DC-OPF. For SE, it is more robust than classical weighted least squares (WLS) to noisy measurements and grid parameters, and it always returns a high-quality estimate even when WLS fails to converge.

Together, the unified architecture and development framework offer a new way to perform large-scale steady-state grid analysis, reduce integration effort, and lower the barrier to entry for power system engineers, marking a key step toward Grid Foundation Models.

\vspace{-5pt}

\end{abstract}